\chapter{Hardware and Data Fundamentals}



This chapter describes the hardware platform upon which the system is built. It covers the dual-sensor radar configuration, the on-chip firmware, and the structure of the data output from each sensor. Understanding the hardware capabilities and output format is essential context for the fusion and detection pipeline described in subsequent chapters. 

\section{Dual-Sensor Hardware Configuration}

The system employs a heterogeneous radar pair from Texas Instruments: one AWR6843AOP (automotive-grade) and one IWR6843AOP (industrial-grade). Both devices are functionally equivalent for this application and share the same physical dimensions, antenna layout, and supported firmware. The use of one unit from each product family reflects hardware availability during development; the fusion pipeline treats both sensors uniformly, with no special handling for minor differences between the two product lines. 

The choice of AWR6843AOP and IWR6843AOP was driven by several considerations: (1) the Antenna-on-Package (AOP) design reduces form factor and eliminates external RF connectors, simplifying physical deployment; (2) both devices integrate all signal processing on-chip, reducing the computational demand on the host edge device; (3) the 60--64 GHz mmWave band provides the sub-centimetre range resolution and Doppler velocity sensitivity required to discriminate between standing, walking, sitting, and lying postures; and (4) the Texas Instruments ecosystem provides a mature firmware ecosystem and open-source host software libraries that accelerate development. 


\subsection{AWR6843AOP}

The AWR6843AOP is a single-chip 60--64 GHz FMCW radar designed for automotivegrade applications. It features an Antenna-on-Package (AOP) design producing an extremely compact footprint (10 \textit{.} 4 \textit{$\times$} 10 \textit{.} 4 mm). The device hosts an on-chip ARM CortexR4F and a dedicated hardware accelerator (HWA) for FFT-based signal processing, along with a C674x DSP for CFAR detection, DBSCAN clustering, and Kalman filter-based person tracking---all executing internally before any data reaches the host computer. 

\subsection{IWR6843AOP}

The IWR6843AOP is the industrial-grade variant sharing the same silicon architecture. It operates over the same 60--64 GHz frequency band with the same MIMO antenna array (3 TX, 4 RX). The primary differences are extended operating temperature ratings and industrial qualification levels; for indoor fall detection at standard room temperature, both sensors are operationally identical. 

\subsection{Shared Hardware Parameters}

Key parameters common to both sensors are summarized in Table 2.1.

\begin{table}[H]
\caption{Key Radar Hardware Parameters}
\label{tab:summary_1}
\centering
\begin{tabularx}{\textwidth}{@{} l X X @{}}
\toprule
\textbf{\textbf{Parameter}} & \textbf{\textbf{Value}} & \textbf{\textbf{Signifcance}} \\ \midrule
Frequency Band & 60--64 GHz & mmWave; privacy-preserving, works in \\
&  & darkness \\
Bandwidth & 3.99 GHz & Range resolution \textit{$\approx$}4 cm theoretical \\
TX Antennas & 3 & Enables 3D angle-of-arrival estimation \\
RX Antennas & 4 & 12 virtual MIMO elements for spatial res- \\
&  & olution \\
On-chip DSP & ARM R4F + C674x & CFAR, DBSCAN, EKF tracking on-chip \\
USB Baud Rate & 921,600 baud & High-speed UART over USB CDC \\
Frame Rate & \textit{$\approx$}10 FPS & 100 ms frame period; adequate for fall de- \\
&  & tection \\
Operational Range & 0.5--20 m & Covers typical room or corridor \\
Point Cloud Output & \textit{x, y, z, v,}snr & 3D position, radial velocity, signal quality \\ \bottomrule
\end{tabularx}
\end{table}





\section{Host Edge Platform: Raspberry Pi}

The host edge computing platform is a Raspberry Pi single-board computer. Both radar sensors connect to the Raspberry Pi via USB CDC serial at 921,600 baud. The Raspberry Pi handles real-time data ingestion, frame synchronization, auto-calibration, point cloud fusion, feature extraction, and file batching. The core pipeline is implemented in Python using NumPy for vectorized point cloud arithmetic and SciPy for SVD computation. No GPU is required on the edge device; all computations complete in under 10 ms per frame on the ARM CPU. 

\section{Firmware: 3D People Tracking}

Both sensors are flashed with Texas Instruments' ``3D People Tracking'' firmware (lab0012). This firmware offloads computational burden to the sensor's on-chip DSP, leaving the host free for multi-sensor fusion. Internally, the firmware executes: 

1. \textbf{CFAR Detection:} Constant False Alarm Rate thresholding identifies radar reflections above the noise floor, converting the range-Doppler map into a sparse set of detected reflection candidates. 

2. \textbf{DBSCAN Clustering:} Groups nearby reflections into coherent clusters representing distinct physical objects. 

3. \textbf{Extended Kalman Filter (EKF) Tracking:} Associates clusters across frames, maintaining stable track IDs ( `tid` ) and providing smooth centroid estimates. 

\section{Understanding Sensor Output: The JSON Frame}

Raw binary telemetry is transmitted at 921,600 baud over USB using TLV (TypeLength-Value) binary encoding. The host parser synchronizes via an 8-byte magic word ( `0x0201040306050807` ), extracts a 32-byte header, and decodes three TLV structures: `TLV\textit{POINT}CLOUD` ( \textit{x, y, z} , Doppler velocity, SNR), `TLV\textit{TRACK}3D` (centroid and bounding box per tracked person), and `TLV\textit{TRACK}INDEX` (mapping from reflection points to their associated track). These are assembled into a structured JSON frame whose key fields are described in Table 2.2. 


\begin{table}[H]
\caption{JSON Frame Fields from 3D People Tracking Firmware}
\label{tab:summary_2}
\centering
\begin{tabularx}{\textwidth}{@{} l X X @{}}
\toprule
\textbf{\textbf{Field}} & \textbf{\textbf{Type}} & \textbf{\textbf{Description}} \\ \midrule
`timestamp` & Float & System time of frame capture (s) \\
`pointCloud` & Array & Variable-length list of detected points (x, y, z, v, snr, track\_id) \\
`targets` & Array & Variable-length list of tracked people (tid, x, y, z, velocity) \\ \bottomrule
\end{tabularx}
\end{table}




\subsection{Point Cloud vs. Targets: Design Choice}

The `targets` array provides a smooth, noise-reduced centroid per person---ideal for geometric calibration. However, it discards all body-shape information. A fall is characterized by a rapid spatial transition from a tall, narrow point cloud (standing) to a wide, low, flat profile (lying)---shape information that exists only in the raw `pointCloud` array. Accordingly, the calibration module uses `targets` centroids for robust spatial alignment, while the fall detector analyzes the full fused `pointCloud` to capture body-shape signatures. 

